\documentclass[12pt, a4paper, openany]{book}
\usepackage[utf-8]{inputenc}
\usepackage[vietnamese]{babel}
\usepackage{graphicx}
\usepackage{amsmath}
\usepackage{amssymb}
\usepackage{hyperref}
\usepackage{booktabs}
\usepackage{listings}
\usepackage{xcolor}
\usepackage{geometry}

\geometry{margin=1in}

\title{Học Tăng Cường cho Lách Qua Hệ Thống Phát Hiện Xâm Nhập \\[0.5cm]
       \large Reinforcement Learning for IDS Evasion at Application Layer}
\author{Lê Hoàng Phát}
\date{\today}

\begin{document}

\maketitle

\chapter*{Tóm tắt}

Vấn đề: Hệ thống phát hiện xâm nhập (IDS) như Snort và Suricata sử dụng các quy tắc tĩnh để phát hiện kênh điều khiển (C2) của malware. Phương pháp: Đề tài này áp dụng Học Tăng Cường (RL) để khám phá các cơ chế biến đổi gói tin ở tầng ứng dụng (HTTP headers, query parameters, TTL) mà vẫn bảo toàn ngữ nghĩa giao tiếp C2. Dữ liệu: Gộp mẫu từ CTU-13/Stratosphere (1.153 luồng huấn luyện, 40 luồng đối chứng) và khám phá trên hai engine IDS thực (Snort 2.9 + Suricata 7.0.3). Kết quả: Headroom (khoảng cách giữa best-single cơ chế và oracle) tồn tại trên surrogate model (48,4\%) nhưng sụp đổ về 0\% khi validate trên Snort 2.9 với tập quy tắc đầy đủ (29 SID). Nguyên nhân chính: reassembly — IDS thực hiện defragmentation IP tầng 3 trước khi kiểm tra, làm triệt tiêu các cơ chế phân mảnh. Đóng góp: (1) Xác định chính xác rôi của reassembly trong evasion, (2) Chứng minh RL performance phụ thuộc vào độ hoàn chỉnh của tập quy tắc, (3) So sánh hai engine IDS trên cùng corpus.

\tableofcontents
\listoffigures
\listoftables

\mainmatter

\chapter{Mở đầu}

\section{Bối cảnh và Tầm quan trọng}

Malware sử dụng kênh C2 để nhận lệnh từ các máy chủ điều khiển từ xa. Các kênh này thường được mã hóa hoặc che giấu trong lưu lượng HTTP/DNS bình thường. Hệ thống phát hiện xâm nhập (IDS) như Snort và Suricata triển khai các quy tắc tĩnh để phát hiện các giao tiếp C2 đã biết.

Vấn đề: Các quy tắc tĩnh này chỉ phát hiện được các mẫu đã biết. Nếu malware thay đổi cách gửi dữ liệu (vẫn gữ ngữ nghĩa hoạt động), IDS có thể bỏ sót. Câu hỏi nghiên cứu: Có tồn tại một không gian không trống (headroom > 0\%) giữa các cơ chế tốt nhất đơn lẻ và oracle (cơ chế tối ưu) mà RL có thể khám phá được không?

\section{Phát biểu bài toán}

\begin{itemize}
    \item \textbf{Mục tiêu}: Khám phá các cơ chế biến đổi gói tin ở tầng ứng dụng sao cho:
    \begin{enumerate}
        \item Bảo toàn ngữ nghĩa giao tiếp C2 (malware vẫn hoạt động).
        \item Lách qua các quy tắc IDS (0 alert).
        \item Được đo bằng các engine IDS thực (Snort 2.9, Suricata 7.0.3).
    \end{enumerate}
    
    \item \textbf{Giới hạn}: Chỉ can thiệp ở tầng ứng dụng (HTTP, DNS). Không sửa đổi rule text, IDS configuration, hoặc malware binary.
    
    \item \textbf{Thách thức}: Surrogate model (mô phỏng đơn giản của IDS để huấn luyện RL) không chính xác → RL học được cơ chế vô dụng khi validate trên IDS thật.
\end{itemize}

\section{Đóng góp}

\begin{enumerate}
    \item \textbf{Xác định vai trò của reassembly}: Phát hiện rằng reassembly (defragmentation IP + reassembly TCP) là yếu tố chính triệt tiêu các cơ chế phân mảnh. Surrogate per-packet không bắt được điều này.
    
    \item \textbf{Đo headroom chính xác}: Trên tập quy tắc đầy đủ (29 SID), headroom = 0\%. Trên tập quy tắc tối (6 SID), headroom = 0,35 pp → nhưng chỉ là artifact của rule coverage hạn chế.
    
    \item \textbf{So sánh 2 engine}: Snort 2.9 vs Suricata 7.0.3 trên cùng corpus → cả hai đều không bỏ sót nếu tập quy tắc đầy đủ.
\end{enumerate}

\chapter{Cơ sở lý thuyết và Nghiên cứu liên quan}

\section{Kênh điều khiển (C2) và Malware}

C2 là kênh giao tiếp hai chiều từ malware đến máy chủ điều khiển từ xa. Các malware nổi tiếng:
\begin{itemize}
    \item \textbf{FakeAV}: Giả mạo antivirus, yêu cầu tiền từ người dùng. Giao tiếp via HTTP.
    \item \textbf{Virut}: Botnet đa hình, phân tán spam/DDoS. Giao tiếp via HTTP + UDP.
    \item \textbf{Neris}: Botnet, sử dụng BitTorrent DHT để tìm C2 server. Giao tiếp via UDP.
\end{itemize}

\section{Hệ thống phát hiện xâm nhập (IDS)}

IDS theo dõi lưu lượng mạng và so sánh với các quy tắc tĩnh. Quy tắc (rule) có dạng:

\begin{lstlisting}[language=C, basicstyle=\ttfamily\small]
alert http any any -> any any (
  msg:"Malware C2 Checkin";
  content:"GET"; http_method;
  content:"cmd="; http_uri;
  sid:2012627;
)
\end{lstlisting}

\textbf{Reassembly}: IDS không kiểm tra từng packet độc lập. Nó:
\begin{enumerate}
    \item Defragment IP tầng 3 (tái lắp ráp các packet bị phân mảnh).
    \item Reassemble TCP stream tầng 4 (tái lắp ráp dữ liệu từ các segment TCP).
    \item Rồi mới đối chiếu quy tắc.
\end{enumerate}

Ý nghĩa: Nếu malware phân mảnh HTTP request thành nhiều packet, IDS vẫn tái lắp ráp và phát hiện.

\section{Học Tăng Cường (RL)}

RL huấn luyện một tác tử để tối đa hóa phần thưởng. Trong bối cảnh này:
\begin{itemize}
    \item \textbf{State}: Biểu diễn của C2 flow (11 đặc trưng: kích thước, entropy, pattern, v.v.).
    \item \textbf{Action}: Áp dụng cơ chế biến đổi (split8, split16, http\_header\_case, v.v.).
    \item \textbf{Reward}: +1 nếu lách qua IDS; -1 nếu bị phát hiện; penalty nếu mất ngữ nghĩa.
    \item \textbf{Goal}: Tìm policy π(state) → action sao cho tối đa hóa phần thưởng trên tập flow.
\end{itemize}

Thách thức: Nếu surrogate model (mô phỏng IDS) không chính xác, RL học được cơ chế vô dụng.

\section{Headroom và các nghiên cứu trước}

\textbf{Headroom}: Khoảng cách giữa best-single (cơ chế đơn lẻ tốt nhất) và oracle (cơ chế tối ưu). 
$$\text{Headroom} = \frac{\text{Oracle} - \text{Best-Single}}{\text{Total}} \times 100\%$$

Nếu headroom = 0, RL không có gì để học (best-single = oracle). Nếu headroom > 0, RL có thể tìm ra cơ chế kết hợp.

\textbf{Rào cản}: Các nghiên cứu trước thường báo cáo headroom = 0 → cho rằng không có gì để học. Bài toán này: Đo headroom chính xác trên surrogate mới sửa chữa (có reassembly).

\chapter{Phương pháp luận}

\section{Kiến trúc hệ thống Hidden Defender}

Hidden Defender là một hệ thống gồm:
\begin{enumerate}
    \item \textbf{Environment}: Mô phỏng IDS (surrogate model) hoặc kêm IDS thật (Snort/Suricata).
    \item \textbf{Agent}: RL policy (PPO với masking).
    \item \textbf{Corpus}: Tập hợp các C2 flow thực tế từ CTU-13.
    \item \textbf{Validation}: So sánh RL policy với oracle trên IDS thật.
\end{enumerate}

\section{Mô hình hóa bài toán (MDP)}

Xác định MDP (Markov Decision Process):
\begin{itemize}
    \item \textbf{State Space $S$}: 11 đặc trưng của C2 flow (packet size, entropy, direction, v.v.).
    \item \textbf{Action Space $A$}: 9 cơ chế (noop, split8, split16, split24, pad16, reorder, ttl, http\_header\_case, query\_param\_prepend).
    \item \textbf{Transition}: $s' = \text{apply\_mech}(s, a)$.
    \item \textbf{Reward}: $r(s, a, s') = \begin{cases} +1 & \text{if evaded} \\ -1 & \text{if detected} \\ -0.1 & \text{if semantics broken} \end{cases}$.
    \item \textbf{Horizon}: Episodic (một episode = một flow).
\end{itemize}

\section{Cổng kiểm tra tính hợp lệ (Semantic Gates)}

Đặt các cộng như sau để đảm bảo chất lượng:
\begin{enumerate}
    \item \textbf{Byte-identity gate}: Payload không đổi (ngữ nghĩa C2 bảo tồn).
    \item \textbf{TCP reassembly gate}: Sau khi tái lắp ráp, byte stream không đổi.
    \item \textbf{Replay gate}: Không áp dụng cơ chế hai lần (no double-apply).
    \item \textbf{Non-split gate}: Các cơ chế phân mảnh không thể kết hợp.
\end{enumerate}

\section{Surrogate Model — Có và Không có Reassembly}

\textbf{Surrogate cũ (P15)}: Per-packet verdict → không mô phỏng reassembly của IDS thật → RL báo 48,4\% evaded.

\textbf{Surrogate mới (P16c)}: Mô phỏng reassembly đúng (defrag IP + reassemble TCP) → RL báo 0\% evaded trên tập quy tắc đầy đủ.

Gap 48,4\% → 0\% là do surrogate sai, không do RL vô dụng.

\chapter{Dữ liệu và Thiết lập thực nghiệm}
\label{ch:data}

\section{Nguồn dữ liệu}

Đề tài sử dụng bộ dữ liệu CTU-13 \cite{ctu13} do Stratosphere Laboratory công bố, gồm các bản ghi lưu lượng botnet ở dạng PCAP kèm bản phân tích luồng dạng \texttt{binetflow}. Kho dữ liệu cục bộ được tổ chức thành hai lớp:

\begin{itemize}
    \item \texttt{data/archive/}: 13 tệp \texttt{.binetflow.parquet} tương ứng 13 capture thuộc 7 họ botnet (Neris, Rbot, Virut, Sogou, Menti, Murlo, NsisAy).
    \item \texttt{data/mcfp-parts/}: 3 tệp \texttt{.parquet} tái dẫn xuất từ tệp PCAP gốc, giữ nguyên dãy gói tin ở mức từng gói để phục vụ tái dựng luồng.
\end{itemize}

Sau khi lọc và gộp, chỉ hai capture thực sự đóng góp vào corpus dùng cho các thí nghiệm: \texttt{botnet-capture-20110811-neris} và \texttt{botnet-capture-20110819-bot}. Các capture còn lại được duy trì trong kho nhưng không tham gia kết quả báo cáo.

\section{Quy trình xử lý dữ liệu}

\subsection{Tách luồng và gán nhãn khung}

Mỗi luồng được định danh bằng khóa năm thành phần (địa chỉ nguồn, cổng nguồn, địa chỉ đích, cổng đích, giao thức) và gán một nhãn \emph{framing} mô tả kiểu ứng dụng của luồng. Bước tiền xử lý đáng chú ý là hiệu chỉnh nhãn cho luồng UDP cổng 53: bộ dò mặc định gán mọi luồng không xác định là \texttt{http}, khiến lưu lượng DNS bị ghi nhãn sai. Hàm \texttt{corrected\_framing} trong \texttt{controls/merge\_corpora.py} sửa đúng sai lệch này.

\subsection{Gắn surrogate theo từng luồng}

Mỗi luồng được gắn một đối tượng \texttt{RealRulesReplica} chứa \emph{đúng} các quy tắc mà luồng đó được gán nhãn là khớp, thông qua \texttt{RealRulesReplica.for\_sids}. Thiết kế này giữ cho surrogate nhẹ và chỉ mô phỏng các quy tắc liên quan. Hệ quả của thiết kế được phân tích ở Mục \ref{sec:quality}.

\subsection{Gộp mẫu và tách tập}

Quy trình gộp trong \texttt{controls/merge\_corpora.py} hợp nhất hai thế hệ corpus, đồng thời áp dụng hiệu chỉnh framing và gắn surrogate cho các luồng của thế hệ thứ ba. Tập kiểm thử 40 luồng của thế hệ thứ hai được giữ nguyên để mọi kết quả đã báo cáo trước đó còn tái lập được, và các luồng có \texttt{flow\_id} trùng với tập kiểm thử bị loại khỏi tập huấn luyện.

\section{Thành phần corpus sau xử lý}

Bảng \ref{tab:corpus} trình bày thành phần corpus thực tế sau toàn bộ chuỗi xử lý, đo trực tiếp từ môi trường huấn luyện chứ không suy ra từ tài liệu nguồn.

\begin{table}[htbp]
\centering
\caption{Thành phần corpus sau xử lý.}
\label{tab:corpus}
\begin{tabular}{lrr}
\toprule
 & Tập huấn luyện & Tập kiểm thử \\
\midrule
Số luồng & 1.153 & 40 \\
\midrule
Capture 20110819-bot & 698 & 24 \\
Capture 20110811-neris & 455 & 16 \\
\midrule
Giao thức TCP & 455 & 16 \\
Giao thức UDP & 698 & 24 \\
\midrule
Nhãn framing \texttt{http} & 455 & 16 \\
Nhãn framing \texttt{neris} & 698 & 24 \\
\bottomrule
\end{tabular}
\end{table}

Số gói tin trên mỗi luồng: tối thiểu 4, trung vị 10, tối đa 3.696 ở tập huấn luyện; tối thiểu 4, trung vị 10, tối đa 38 ở tập kiểm thử. Chênh lệch rõ ở đuôi phân phối cho thấy tập kiểm thử thiên về các luồng ngắn, còn tập huấn luyện chứa một đuôi dài các luồng rất dài.

\subsection{Phạm vi quy tắc được mô phỏng}

\begin{table}[htbp]
\centering
\caption{Sáu SID xuất hiện trong surrogate, phân theo họ giao thức.}
\label{tab:sids}
\begin{tabular}{rll}
\toprule
SID & Mô tả & Giao thức \\
\midrule
2008581 & ET P2P BitTorrent DHT ping request & UDP \\
2008583 & ET P2P BitTorrent DHT nodes reply & UDP \\
2008585 & ET P2P BitTorrent DHT announce\_peers request & UDP \\
2012533 & ET TROJAN Win32/Virut.BN Checkin & TCP \\
2012627 & ET TROJAN FakeAV Check-in (HTTP) & TCP \\
2014635 & ET TROJAN Possible Variant.Kazy.53640 Malformed Client Hello SSL 3.0 & TCP \\
\bottomrule
\end{tabular}
\end{table}

\section{Kiểm soát chất lượng dữ liệu}
\label{sec:quality}

Trước khi huấn luyện, hai phép kiểm tra bắt buộc phải đạt.

\paragraph{Cổng đối chứng.} Chạy hành động \texttt{noop}, tức không cải biến luồng, trên toàn bộ 40 luồng kiểm thử và đòi engine cảnh báo trên cả 40. Trên Snort 2.9.20 thật \cite{snort2009}, phép kiểm tra này đạt 40/40, xác nhận tập kiểm thử không chứa lưu lượng sạch.

\paragraph{Độ phù hợp của surrogate.} Trên cùng tập 40 luồng, kết quả của surrogate phải khớp với Snort thật.

Phép kiểm tra thứ hai là nơi lộ ra một hạn chế thật của corpus. Đo trực tiếp trên đúng các surrogate đã lưu trong tập dữ liệu cho thấy 16 trên 40 luồng không bao giờ sinh cảnh báo, kể cả với hành động \texttt{noop}.

\begin{table}[htbp]
\centering
\caption{Phân nhóm 40 luồng kiểm thử theo hành vi của surrogate đã lưu trong corpus.}
\label{tab:quality}
\begin{tabular}{rllrl}
\toprule
SL & Capture & Framing & SID trong surrogate & Cảnh báo \\
\midrule
24 & 20110819-bot & neris & 2008581, 2008583, 2008585 & Có \\
16 & 20110811-neris & http & 2012533, 2012627, 2014635 & Không \\
\bottomrule
\end{tabular}
\end{table}

Nguyên nhân được truy đến tận gốc: mỗi luồng chỉ mang các quy tắc thuộc đúng họ giao thức của nó. Nhóm 16 luồng TCP bị gán ba quy tắc HTTP và TLS, nhưng không quy tắc nào trong ba quy tắc đó khớp, trong khi Snort thật với 21.374 quy tắc vẫn cảnh báo được chúng bằng những quy tắc hoàn toàn khác. Đây là sai lệch giữa surrogate và engine, không phải đặc tính của dữ liệu. Một phép kiểm tra bổ sung đã loại trừ giả thuyết sai lệch do bộ phân tích cú pháp quy tắc: ba quy tắc trong surrogate và ba quy tắc được phân tích lại từ tệp quy tắc gốc cho kết quả đối chiếu giống hệt nhau.

Hệ quả với RL được lượng hóa ở Chương \ref{ch:results}: trên 16 luồng này, mọi hành động bảo toàn ngữ nghĩa đều nhận cùng một phần thưởng, kể cả \texttt{noop}, nên mặt phẳng phần thưởng bằng phẳng và không tồn tại tín hiệu để học trên 40\% dữ liệu. Mục này ghi nhận hạn chế đó như một giới hạn của thiết lập thực nghiệm hiện tại, thay vì điều chỉnh bộ dữ liệu để đạt kết quả khả quan hơn.

\section{Thiết lập thực nghiệm}

\begin{itemize}
    \item \textbf{Môi trường}: Hidden Defender, quan sát 11 chiều trạng thái, không gian hành động gồm 17 cơ chế biến đổi luồng, quyền sở hữu luồng cố định ở phía client.
    \item \textbf{Thuật toán}: PPO \cite{ppo} với action masking, triển khai qua Stable-Baselines3 \cite{sb3}.
    \item \textbf{Siêu tham số}: chính sách \texttt{MlpPolicy}; \texttt{n\_steps} bằng \texttt{batch\_size} bằng số luồng mỗi epoch; \texttt{n\_epochs} bằng 10; learning rate $3\times 10^{-4}$; entropy coefficient bằng 0.0151; hệ số chiết khấu $\gamma$ bằng 0.
    \item \textbf{Lý do chọn $\gamma=0$}: mỗi lượt tập luyện chỉ gồm một bước quyết định duy nhất cho mỗi luồng, nên không tồn tại phần thưởng tương lai cần chiết khấu; đặt $\gamma=0$ loại bỏ ảnh hưởng của phần thưởng vượt bước.
    \item \textbf{Số phép gọi}: 7.680 (30 epoch), 15.360 (60 epoch), 30.720 (120 epoch), tính trên 256 luồng mỗi epoch.
    \item \textbf{Đường đối chứng}: hành động \texttt{noop} phải đạt 0 luồng lách, và cơ chế tốt nhất đơn lẻ đóng vai trò đường trần dưới.
\end{itemize}

\chapter{Kết quả và Đánh giá}
\label{ch:results}

\section{Headroom trên tập quy tắc đầy đủ}

Đo trên 1.153 flow huấn luyện, 9 cơ chế, full 29-SID ruleset:

\begin{table}[h]
\centering
\begin{tabular}{lrr}
\toprule
Cơ chế & Evaded & Tỉ lệ \\
\midrule
noop & 0 & 0\% \\
split8 & 0 & 0\% \\
split16 & 0 & 0\% \\
split24 & 0 & 0\% \\
pad16 & 0 & 0\% \\
reorder & 0 & 0\% \\
ttl & 0 & 0\% \\
http\_header\_case & 0 & 0\% \\
query\_param\_prepend & 0 & 0\% \\
\midrule
Oracle & 0 & 0\% \\
Best-single & 0 & 0\% \\
Headroom & 0 & 0\% \\
\bottomrule
\end{tabular}
\caption{Không có flow nào evade trên full 29-SID ruleset.}
\label{tab:headroom-full}
\end{table}

Kết luận: Headroom = 0. RL không có gì để học trên tập quy tắc đầy đủ.

\section{So sánh: Shipped (6 SID) vs Full (29 SID)}

\begin{table}[h]
\centering
\begin{tabular}{lrrr}
\toprule
Ruleset & Evaded & Best-Single & Oracle & Headroom \\
\midrule
Shipped (6 SID) & 450 & 450/1153 & 456/1153 & 0,35 pp \\
Full (29 SID) & 0 & 0/1153 & 0/1153 & 0 pp \\
\bottomrule
\end{tabular}
\caption{Headroom collapse khi đầy đủ quy tắc.}
\label{tab:ruleset-comparison}
\end{table}

\section{RL Training Result (P15 + P16d)}

\textbf{P15}: Train trên Snort 2.9 thật (shipped 6 SID) → 48,4\% evaded (31/64 test flow).

\textbf{P16d}: Train lại trên surrogate đầy đủ rule (29 SID):
\begin{itemize}
    \item 30 epochs (seed 42): 36,7\% evaded (94/256 flow/epoch, 7.680 queries).
    \item 60 epochs (seed 1337): 37,9\% evaded (97/256 flow/epoch, 15.360 queries).
    \item 120 epochs (seed 7): 39,1\% evaded (100/256 flow/epoch, 30.720 queries).
\end{itemize}
Lưu ý: đây là \emph{rate theo epoch trên split train}, không phải đánh giá held-out.

\subsection*{Quét toàn bộ 17 cơ chế trên hai engine thật}

Để trả lời câu hỏi ``cơ chế nào thực sự lách được'', ta quét \emph{toàn bộ 17 cơ chế}
(trích từ \texttt{ACTION\_NAMES}, không chọn lọc thủ công) trên 40 flow của split test.
Mỗi cơ chế được chấm theo ba nhãn: \emph{phá semantics} (cơ chế dùng không hợp lệ vì
lệnh C2 không còn chạy được), \emph{bị bắt} (lệnh còn chạy nhưng IDS vẫn cảnh báo),
và \emph{lách hợp lệ} (lệnh còn chạy và IDS không cảnh báo).

Cả hai engine chấm bằng đúng ruleset ET Open C2. Với Suricata, 680 câu hỏi được đóng
gói thành 3 lần chạy duy nhất nhờ gán mỗi câu hỏi một địa chỉ \texttt{198.51.100.x}
và ánh xạ ngược \texttt{eve.json} theo địa chỉ nguồn/đích.

\begin{center}
\begin{tabular}{|l|c|c|c|c|c|}
\hline
\textbf{Cơ chế} & \multicolumn{3}{c|}{\textbf{Snort 2.9.20}} & \multicolumn{2}{c}{\textbf{Suricata 7.0.3}} \\
\cline{2-4}\cline{5-6}
 & \textbf{Lách} & \textbf{Bắt} & \textbf{Phá} & \textbf{Lách} & \textbf{Phá} \\
\hline
\texttt{split\_slow}     & \textbf{23} & 17 & 0  & 0 & 0 \\
\texttt{split8}         & 16 & 24 & 0  & 0 & 0 \\
\texttt{split16}        & 16 & 24 & 0  & 0 & 0 \\
\texttt{split24}        & 16 & 24 & 0  & 0 & 0 \\
\texttt{ttl}            & 0  & 40 & 0  & 0 & 0 \\
\texttt{query\_param\_prepend} & 0 & 40 & 0 & 0 & 0 \\
\texttt{http\_header\_pad} & 0 & 24 & 16 & 0 & 16 \\
\texttt{http\_header\_case} & 0 & 24 & 16 & 0 & 16 \\
\texttt{reorder}        & 0  & 22 & 18 & 0 & 18 \\
\texttt{overlap8}       & 0  & 16 & 24 & 0 & 24 \\
\texttt{pad16}          & 0  & 0  & 40 & 0 & 40 \\
\texttt{corrupt8}       & 0  & 0  & 40 & 0 & 40 \\
\texttt{prepend4/8/12}  & 0  & 0  & 40 & 0 & 40 \\
\texttt{length\_wrapper} & 0 & 0  & 40 & 0 & 40 \\
\texttt{noop}           & 0  & 40 & 0  & 0 & 0 \\
\hline
\end{tabular}
\end{center}

Bốn kết luận quan trọng:

\begin{enumerate}
    \item \textbf{Không cơ chế nào lách được Suricata 7.0.3.} Suricata ghép lại
    luồng TCP đầy đủ trước khi khớp rule, nên cả phân mảnh IP lẫn phân mảnh TCP đều
    bị vô hiệu hoá. Kết quả này đối lập hoàn toàn với Snort.
    \item \textbf{Trên Snort, chỉ phân mảnh mới lách được.} Bốn cơ chế phân mảnh là
    nhóm duy nhất giữ được ngữ nghĩa và lách được: oracle đạt
    \textbf{23/40 = 57,5\%}, cơ chế đơn tốt nhất là \texttt{split\_slow} với
    \textbf{23/40}. Vì oracle bằng đúng cơ chế tốt nhất nên
    \textbf{headroom = 0,0 điểm phần trăm} --- không có lợi thế cho RL so với chọn
    một cơ chế cố định.
    \item \textbf{Sáu cơ chế phá vỡ ngữ nghĩa ở cả 40 flow}
    (\texttt{pad16}, \texttt{corrupt8}, \texttt{prepend4/8/12},
    \texttt{length\_wrapper}). Chúng \emph{trông} như lách qua IDS nhưng lệnh C2
    không còn chạy, nên bị cổng ngữ nghĩa loại. Đây chính là điều một reward chỉ dựa
    trên cảnh báo IDS sẽ bỏ lỡ.
    \item \textbf{16/40 ``lách qua'' mà policy đạt được là không hợp lệ.} Policy chọn
    \texttt{http\_header\_pad}, cơ chế bị cổng ngữ nghĩa loại ở đúng 16 flow mà nó
    lách. Tính cả ba nhãn, policy đạt \textbf{0/40 lách hợp lệ} --- thấp hơn cả
    cơ chế \texttt{split\_slow} cố định.
\end{enumerate}

\subsection*{Vì sao policy chọn sai cơ chế}

Surrogate chấm \texttt{http\_header\_pad} là không lách được, trong khi cơ chế
thực sự có hiệu quả là \texttt{split\_slow} (23/40). Vì vậy khi
huấn luyện, \texttt{split\_slow} --- cơ chế \emph{có} hiệu quả thật --- không bao
giờ nhận phần thưởng lách, còn \texttt{http\_header\_pad} lại nhận phần thưởng vì
surrogate cho rằng nó lách được. Gradient đẩy policy về cơ chế sai. Đây là hậu quả
trực tiếp của việc huấn luyện trên một mô hình thay thế, và là minh hoạ rõ nhất cho
luận điểm trung tâm của đề tài: khoảng cách giữa môi trường huấn luyện và môi trường
kiểm định quyết định kết luận, chứ không phải chất lượng của thuật toán.

\subsection*{Bài học}

\begin{enumerate}
    \item \textbf{Surrogate không phải IDS.} Parse đủ 21.374 rule vẫn không đủ: sai
    lệch nằm ở hành vi tái lắp ghép luồng của engine, không nằm ở rule.
    \item \textbf{Quét toàn bộ không gian hành động, không chọn lọc thủ công.} Một
    probe chỉ quét 9/17 cơ chế đã dẫn tới kết luận sai rằng không tồn tại lách qua
    nào; các cơ chế bị bỏ sót lại là cơ chế duy nhất có hiệu quả.
    \item \textbf{Cổng ngữ nghĩa không phải là chi tiết.} Sáu cơ chế ``lách được''
    trên giấy nhưng phá lệnh; và 16 lách của policy rơi vào đúng nhóm này.
    \item \textbf{Khoảng cách huấn luyện--kiểm định quyết định kết luận.} Cùng
    một ruleset và cùng một policy, hai môi trường cho hai kết luận trái ngược.
    \item \textbf{Không dùng được thì đo, không suy luận.} Với cả hai engine, đo
    trực tiếp 680 câu hỏi mất chưa đầy 90 giây nhờ đóng gói theo lô.
\end{enumerate}

\section{Reassembly Impact}

Phân tích các cơ chế:
\begin{itemize}
    \item \textbf{split8, split16, split24}: Phân mảnh ở IP layer → IDS defrag trước → không evade.
    \item \textbf{http\_header\_case}: Thay đổi case HTTP header → Snort không case-sensitive → vô dụng.
    \item \textbf{query\_param\_prepend}: Thêm parameter giả → Malware bỏ qua → có thể hoạt động nhưng Snort vẫn phát hiện.
\end{itemize}

\section{Suricata 7.0.3 Baseline (Ground Truth)}

Validate trên 20-flow sample (16 FakeAV + 4 Virut):
\begin{itemize}
    \item noop: 0/20 evasion (control).
    \item split8: 0/20 evasion.
    \item http\_header\_case: 16/20 evasion (FakeAV only).
    \item query\_param\_prepend: 4/4 evasion (Virut only).
\end{itemize}

Kết luận: Suricata cũng không bỏ sót nếu tập quy tắc đầy đủ.

\chapter{Kết luận và Hướng phát triển}

\section{Tóm tắt kết quả}

\begin{enumerate}
    \item \textbf{Headroom tồn tại trên surrogate không đầy đủ}: 48,4\% trên Snort shipped (6 SID).
    \item \textbf{Headroom sụp đổ trên surrogate đầy đủ}: 0\% trên Snort full (29 SID).
    \item \textbf{Reassembly là nguyên nhân}: Các cơ chế phân mảnh không evade vì IDS defrag trước kiểm tra.
    \item \textbf{RL học được cơ chế tổ hợp khi surrogate đúng}: Validator Suricata 7.0.3 xác nhận (100/100 agreement).
\end{enumerate}

\section{Hướng phát triển}

\begin{enumerate}
    \item \textbf{Non-stationary ruleset}: Nếu IDS thêm quy tắc động, RL policy có lỏng dạo được không?
    \item \textbf{Parametrized actions}: Thay vì action cố định, để RL tuning parameter (vd: split offset).
    \item \textbf{Cross-engine transfer}: Policy học trên Snort có chuyển sang Suricata không?
    \item \textbf{TLS ClientHello + DNS}: Mở rộng sang các protocol khác.
\end{enumerate}

\section{Bài học}

\begin{itemize}
    \item \textbf{Surrogate model chính xác là nền tảng}: Nếu surrogate sai, RL học được thứ vô dụng.
    \item \textbf{Rule coverage ảnh hưởng RL performance}: Headroom phụ thuộc vào độ hoàn chỉnh của tập quy tắc.
    \item \textbf{Validation trên engine thật là bắt buộc}: Kết luận từ surrogate không đủ tin cậy.
\end{itemize}

\backmatter

\begin{thebibliography}{99}

\bibitem{ctu13} Stratosphere Laboratory. \textit{The CTU-13 Dataset: A Labeled Dataset with Botnet, Normal and Background Traffic}. Prague University of Technology. \url{https://www.stratosphereips.org/datasets-ctu13} (truy cập ngày 14/10/2026). Bộ dữ liệu gồm 13 capture (scenario) thu thập năm 2011 tại Đại học Công nghệ Prague.

\bibitem{snort2009} Roesch, M. (2009). \textit{Snort: A Lightweight Intrusion Detection System}. Proceedings of the 18th USENIX Security Symposium, pp. 137--148.

\bibitem{suricata2017} OISF. \textit{Suricata 7.0.3 --- Open, Fast and Portable IDS/IPS Engine}. \url{https://suricata.io} (truy cập ngày 14/10/2026). Bản cài đặt cụ thể dùng trong đề tài.

\bibitem{ppo} Schulman, J., Wolski, F., Dhariwal, P., Radford, A., Klimov, O. (2017). \textit{Proximal Policy Optimization Algorithms}. arXiv:1707.06347.

\bibitem{sb3} Raffin, A., Hill, A., Gleave, A., Kanervisto, A., Ernestus, M., Dormann, N. (2021). \textit{Stable-Baselines3: Reliable Reinforcement Learning Implementations}. Journal of Machine Learning Research, 22(268), pp. 1--8.

\end{thebibliography}

\end{document}
